\documentclass[11pt]{article}

    \usepackage[breakable]{tcolorbox}
    \usepackage{parskip} % Stop auto-indenting (to mimic markdown behaviour)
    

    % Basic figure setup, for now with no caption control since it's done
    % automatically by Pandoc (which extracts ![](path) syntax from Markdown).
    \usepackage{graphicx}
    % Keep aspect ratio if custom image width or height is specified
    \setkeys{Gin}{keepaspectratio}
    % Maintain compatibility with old templates. Remove in nbconvert 6.0
    \let\Oldincludegraphics\includegraphics
    % Ensure that by default, figures have no caption (until we provide a
    % proper Figure object with a Caption API and a way to capture that
    % in the conversion process - todo).
    \usepackage{caption}
    \DeclareCaptionFormat{nocaption}{}
    \captionsetup{format=nocaption,aboveskip=0pt,belowskip=0pt}

    \usepackage{float}
    \floatplacement{figure}{H} % forces figures to be placed at the correct location
    \usepackage{xcolor} % Allow colors to be defined
    \usepackage{enumerate} % Needed for markdown enumerations to work
    \usepackage{geometry} % Used to adjust the document margins
    \usepackage{amsmath} % Equations
    \usepackage{amssymb} % Equations
    \usepackage{textcomp} % defines textquotesingle
    % Hack from http://tex.stackexchange.com/a/47451/13684:
    \AtBeginDocument{%
        \def\PYZsq{\textquotesingle}% Upright quotes in Pygmentized code
    }
    \usepackage{upquote} % Upright quotes for verbatim code
    \usepackage{eurosym} % defines \euro

    \usepackage{iftex}
    \ifPDFTeX
        \usepackage[T1]{fontenc}
        \IfFileExists{alphabeta.sty}{
              \usepackage{alphabeta}
          }{
              \usepackage[mathletters]{ucs}
              \usepackage[utf8x]{inputenc}
          }
    \else
        \usepackage{fontspec}
        \usepackage{unicode-math}
    \fi

    \usepackage{fancyvrb} % verbatim replacement that allows latex
    \setcounter{secnumdepth}{0}
    \usepackage{grffile} % extends the file name processing of package graphics
                         % to support a larger range
    \makeatletter % fix for old versions of grffile with XeLaTeX
    \@ifpackagelater{grffile}{2019/11/01}
    {
      % Do nothing on new versions
    }
    {
      \def\Gread@@xetex#1{%
        \IfFileExists{"\Gin@base".bb}%
        {\Gread@eps{\Gin@base.bb}}%
        {\Gread@@xetex@aux#1}%
      }
    }
    \makeatother
    \usepackage[Export]{adjustbox} % Used to constrain images to a maximum size
    \adjustboxset{max size={0.9\linewidth}{0.9\paperheight}}

    % The hyperref package gives us a pdf with properly built
    % internal navigation ('pdf bookmarks' for the table of contents,
    % internal cross-reference links, web links for URLs, etc.)
    \usepackage{hyperref}
    % The default LaTeX title has an obnoxious amount of whitespace. By default,
    % titling removes some of it. It also provides customization options.
    \usepackage{titling}
    \usepackage{longtable} % longtable support required by pandoc >1.10
    \usepackage{booktabs}  % table support for pandoc > 1.12.2
    \usepackage{array}     % table support for pandoc >= 2.11.3
    \usepackage{calc}      % table minipage width calculation for pandoc >= 2.11.1
    \usepackage[inline]{enumitem} % IRkernel/repr support (it uses the enumerate* environment)
    \usepackage[normalem]{ulem} % ulem is needed to support strikethroughs (\sout)
                                % normalem makes italics be italics, not underlines
    \usepackage{soul}      % strikethrough (\st) support for pandoc >= 3.0.0
    \usepackage{mathrsfs}
    

    
    % Colors for the hyperref package
    \definecolor{urlcolor}{rgb}{0,.145,.698}
    \definecolor{linkcolor}{rgb}{.71,0.21,0.01}
    \definecolor{citecolor}{rgb}{.12,.54,.11}

    % ANSI colors
    \definecolor{ansi-black}{HTML}{3E424D}
    \definecolor{ansi-black-intense}{HTML}{282C36}
    \definecolor{ansi-red}{HTML}{E75C58}
    \definecolor{ansi-red-intense}{HTML}{B22B31}
    \definecolor{ansi-green}{HTML}{00A250}
    \definecolor{ansi-green-intense}{HTML}{007427}
    \definecolor{ansi-yellow}{HTML}{DDB62B}
    \definecolor{ansi-yellow-intense}{HTML}{B27D12}
    \definecolor{ansi-blue}{HTML}{208FFB}
    \definecolor{ansi-blue-intense}{HTML}{0065CA}
    \definecolor{ansi-magenta}{HTML}{D160C4}
    \definecolor{ansi-magenta-intense}{HTML}{A03196}
    \definecolor{ansi-cyan}{HTML}{60C6C8}
    \definecolor{ansi-cyan-intense}{HTML}{258F8F}
    \definecolor{ansi-white}{HTML}{C5C1B4}
    \definecolor{ansi-white-intense}{HTML}{A1A6B2}
    \definecolor{ansi-default-inverse-fg}{HTML}{FFFFFF}
    \definecolor{ansi-default-inverse-bg}{HTML}{000000}

    % common color for the border for error outputs.
    \definecolor{outerrorbackground}{HTML}{FFDFDF}

    % commands and environments needed by pandoc snippets
    % extracted from the output of `pandoc -s`
    \providecommand{\tightlist}{%
      \setlength{\itemsep}{0pt}\setlength{\parskip}{0pt}}
    \DefineVerbatimEnvironment{Highlighting}{Verbatim}{commandchars=\\\{\}}
    % Add ',fontsize=\small' for more characters per line
    \newenvironment{Shaded}{}{}
    \newcommand{\KeywordTok}[1]{\textcolor[rgb]{0.00,0.44,0.13}{\textbf{{#1}}}}
    \newcommand{\DataTypeTok}[1]{\textcolor[rgb]{0.56,0.13,0.00}{{#1}}}
    \newcommand{\DecValTok}[1]{\textcolor[rgb]{0.25,0.63,0.44}{{#1}}}
    \newcommand{\BaseNTok}[1]{\textcolor[rgb]{0.25,0.63,0.44}{{#1}}}
    \newcommand{\FloatTok}[1]{\textcolor[rgb]{0.25,0.63,0.44}{{#1}}}
    \newcommand{\CharTok}[1]{\textcolor[rgb]{0.25,0.44,0.63}{{#1}}}
    \newcommand{\StringTok}[1]{\textcolor[rgb]{0.25,0.44,0.63}{{#1}}}
    \newcommand{\CommentTok}[1]{\textcolor[rgb]{0.38,0.63,0.69}{\textit{{#1}}}}
    \newcommand{\OtherTok}[1]{\textcolor[rgb]{0.00,0.44,0.13}{{#1}}}
    \newcommand{\AlertTok}[1]{\textcolor[rgb]{1.00,0.00,0.00}{\textbf{{#1}}}}
    \newcommand{\FunctionTok}[1]{\textcolor[rgb]{0.02,0.16,0.49}{{#1}}}
    \newcommand{\RegionMarkerTok}[1]{{#1}}
    \newcommand{\ErrorTok}[1]{\textcolor[rgb]{1.00,0.00,0.00}{\textbf{{#1}}}}
    \newcommand{\NormalTok}[1]{{#1}}

    % Additional commands for more recent versions of Pandoc
    \newcommand{\ConstantTok}[1]{\textcolor[rgb]{0.53,0.00,0.00}{{#1}}}
    \newcommand{\SpecialCharTok}[1]{\textcolor[rgb]{0.25,0.44,0.63}{{#1}}}
    \newcommand{\VerbatimStringTok}[1]{\textcolor[rgb]{0.25,0.44,0.63}{{#1}}}
    \newcommand{\SpecialStringTok}[1]{\textcolor[rgb]{0.73,0.40,0.53}{{#1}}}
    \newcommand{\ImportTok}[1]{{#1}}
    \newcommand{\DocumentationTok}[1]{\textcolor[rgb]{0.73,0.13,0.13}{\textit{{#1}}}}
    \newcommand{\AnnotationTok}[1]{\textcolor[rgb]{0.38,0.63,0.69}{\textbf{\textit{{#1}}}}}
    \newcommand{\CommentVarTok}[1]{\textcolor[rgb]{0.38,0.63,0.69}{\textbf{\textit{{#1}}}}}
    \newcommand{\VariableTok}[1]{\textcolor[rgb]{0.10,0.09,0.49}{{#1}}}
    \newcommand{\ControlFlowTok}[1]{\textcolor[rgb]{0.00,0.44,0.13}{\textbf{{#1}}}}
    \newcommand{\OperatorTok}[1]{\textcolor[rgb]{0.40,0.40,0.40}{{#1}}}
    \newcommand{\BuiltInTok}[1]{{#1}}
    \newcommand{\ExtensionTok}[1]{{#1}}
    \newcommand{\PreprocessorTok}[1]{\textcolor[rgb]{0.74,0.48,0.00}{{#1}}}
    \newcommand{\AttributeTok}[1]{\textcolor[rgb]{0.49,0.56,0.16}{{#1}}}
    \newcommand{\InformationTok}[1]{\textcolor[rgb]{0.38,0.63,0.69}{\textbf{\textit{{#1}}}}}
    \newcommand{\WarningTok}[1]{\textcolor[rgb]{0.38,0.63,0.69}{\textbf{\textit{{#1}}}}}
    \makeatletter
    \newsavebox\pandoc@box
    \newcommand*\pandocbounded[1]{%
      \sbox\pandoc@box{#1}%
      % scaling factors for width and height
      \Gscale@div\@tempa\textheight{\dimexpr\ht\pandoc@box+\dp\pandoc@box\relax}%
      \Gscale@div\@tempb\linewidth{\wd\pandoc@box}%
      % select the smaller of both
      \ifdim\@tempb\p@<\@tempa\p@
        \let\@tempa\@tempb
      \fi
      % scaling accordingly (\@tempa < 1)
      \ifdim\@tempa\p@<\p@
        \scalebox{\@tempa}{\usebox\pandoc@box}%
      % scaling not needed, use as it is
      \else
        \usebox{\pandoc@box}%
      \fi
    }
    \makeatother

    % Define a nice break command that doesn't care if a line doesn't already
    % exist.
    \def\br{\hspace*{\fill} \\* }
    % Math Jax compatibility definitions
    \def\gt{>}
    \def\lt{<}
    \let\Oldtex\TeX
    \let\Oldlatex\LaTeX
    \renewcommand{\TeX}{\textrm{\Oldtex}}
    \renewcommand{\LaTeX}{\textrm{\Oldlatex}}
    % Document parameters
    % Document title
    \title{Analisi dei Segnali del Giroscopio}
    
    
    \date{}
    
    
    
    
    \author{}
    
    
    
% Pygments definitions
\makeatletter
\def\PY@reset{\let\PY@it=\relax \let\PY@bf=\relax%
    \let\PY@ul=\relax \let\PY@tc=\relax%
    \let\PY@bc=\relax \let\PY@ff=\relax}
\def\PY@tok#1{\csname PY@tok@#1\endcsname}
\def\PY@toks#1+{\ifx\relax#1\empty\else%
    \PY@tok{#1}\expandafter\PY@toks\fi}
\def\PY@do#1{\PY@bc{\PY@tc{\PY@ul{%
    \PY@it{\PY@bf{\PY@ff{#1}}}}}}}
\def\PY#1#2{\PY@reset\PY@toks#1+\relax+\PY@do{#2}}

\@namedef{PY@tok@w}{\def\PY@tc##1{\textcolor[rgb]{0.73,0.73,0.73}{##1}}}
\@namedef{PY@tok@c}{\let\PY@it=\textit\def\PY@tc##1{\textcolor[rgb]{0.24,0.48,0.48}{##1}}}
\@namedef{PY@tok@cp}{\def\PY@tc##1{\textcolor[rgb]{0.61,0.40,0.00}{##1}}}
\@namedef{PY@tok@k}{\let\PY@bf=\textbf\def\PY@tc##1{\textcolor[rgb]{0.00,0.50,0.00}{##1}}}
\@namedef{PY@tok@kp}{\def\PY@tc##1{\textcolor[rgb]{0.00,0.50,0.00}{##1}}}
\@namedef{PY@tok@kt}{\def\PY@tc##1{\textcolor[rgb]{0.69,0.00,0.25}{##1}}}
\@namedef{PY@tok@o}{\def\PY@tc##1{\textcolor[rgb]{0.40,0.40,0.40}{##1}}}
\@namedef{PY@tok@ow}{\let\PY@bf=\textbf\def\PY@tc##1{\textcolor[rgb]{0.67,0.13,1.00}{##1}}}
\@namedef{PY@tok@nb}{\def\PY@tc##1{\textcolor[rgb]{0.00,0.50,0.00}{##1}}}
\@namedef{PY@tok@nf}{\def\PY@tc##1{\textcolor[rgb]{0.00,0.00,1.00}{##1}}}
\@namedef{PY@tok@nc}{\let\PY@bf=\textbf\def\PY@tc##1{\textcolor[rgb]{0.00,0.00,1.00}{##1}}}
\@namedef{PY@tok@nn}{\let\PY@bf=\textbf\def\PY@tc##1{\textcolor[rgb]{0.00,0.00,1.00}{##1}}}
\@namedef{PY@tok@ne}{\let\PY@bf=\textbf\def\PY@tc##1{\textcolor[rgb]{0.80,0.25,0.22}{##1}}}
\@namedef{PY@tok@nv}{\def\PY@tc##1{\textcolor[rgb]{0.10,0.09,0.49}{##1}}}
\@namedef{PY@tok@no}{\def\PY@tc##1{\textcolor[rgb]{0.53,0.00,0.00}{##1}}}
\@namedef{PY@tok@nl}{\def\PY@tc##1{\textcolor[rgb]{0.46,0.46,0.00}{##1}}}
\@namedef{PY@tok@ni}{\let\PY@bf=\textbf\def\PY@tc##1{\textcolor[rgb]{0.44,0.44,0.44}{##1}}}
\@namedef{PY@tok@na}{\def\PY@tc##1{\textcolor[rgb]{0.41,0.47,0.13}{##1}}}
\@namedef{PY@tok@nt}{\let\PY@bf=\textbf\def\PY@tc##1{\textcolor[rgb]{0.00,0.50,0.00}{##1}}}
\@namedef{PY@tok@nd}{\def\PY@tc##1{\textcolor[rgb]{0.67,0.13,1.00}{##1}}}
\@namedef{PY@tok@s}{\def\PY@tc##1{\textcolor[rgb]{0.73,0.13,0.13}{##1}}}
\@namedef{PY@tok@sd}{\let\PY@it=\textit\def\PY@tc##1{\textcolor[rgb]{0.73,0.13,0.13}{##1}}}
\@namedef{PY@tok@si}{\let\PY@bf=\textbf\def\PY@tc##1{\textcolor[rgb]{0.64,0.35,0.47}{##1}}}
\@namedef{PY@tok@se}{\let\PY@bf=\textbf\def\PY@tc##1{\textcolor[rgb]{0.67,0.36,0.12}{##1}}}
\@namedef{PY@tok@sr}{\def\PY@tc##1{\textcolor[rgb]{0.64,0.35,0.47}{##1}}}
\@namedef{PY@tok@ss}{\def\PY@tc##1{\textcolor[rgb]{0.10,0.09,0.49}{##1}}}
\@namedef{PY@tok@sx}{\def\PY@tc##1{\textcolor[rgb]{0.00,0.50,0.00}{##1}}}
\@namedef{PY@tok@m}{\def\PY@tc##1{\textcolor[rgb]{0.40,0.40,0.40}{##1}}}
\@namedef{PY@tok@gh}{\let\PY@bf=\textbf\def\PY@tc##1{\textcolor[rgb]{0.00,0.00,0.50}{##1}}}
\@namedef{PY@tok@gu}{\let\PY@bf=\textbf\def\PY@tc##1{\textcolor[rgb]{0.50,0.00,0.50}{##1}}}
\@namedef{PY@tok@gd}{\def\PY@tc##1{\textcolor[rgb]{0.63,0.00,0.00}{##1}}}
\@namedef{PY@tok@gi}{\def\PY@tc##1{\textcolor[rgb]{0.00,0.52,0.00}{##1}}}
\@namedef{PY@tok@gr}{\def\PY@tc##1{\textcolor[rgb]{0.89,0.00,0.00}{##1}}}
\@namedef{PY@tok@ge}{\let\PY@it=\textit}
\@namedef{PY@tok@gs}{\let\PY@bf=\textbf}
\@namedef{PY@tok@ges}{\let\PY@bf=\textbf\let\PY@it=\textit}
\@namedef{PY@tok@gp}{\let\PY@bf=\textbf\def\PY@tc##1{\textcolor[rgb]{0.00,0.00,0.50}{##1}}}
\@namedef{PY@tok@go}{\def\PY@tc##1{\textcolor[rgb]{0.44,0.44,0.44}{##1}}}
\@namedef{PY@tok@gt}{\def\PY@tc##1{\textcolor[rgb]{0.00,0.27,0.87}{##1}}}
\@namedef{PY@tok@err}{\def\PY@bc##1{{\setlength{\fboxsep}{\string -\fboxrule}\fcolorbox[rgb]{1.00,0.00,0.00}{1,1,1}{\strut ##1}}}}
\@namedef{PY@tok@kc}{\let\PY@bf=\textbf\def\PY@tc##1{\textcolor[rgb]{0.00,0.50,0.00}{##1}}}
\@namedef{PY@tok@kd}{\let\PY@bf=\textbf\def\PY@tc##1{\textcolor[rgb]{0.00,0.50,0.00}{##1}}}
\@namedef{PY@tok@kn}{\let\PY@bf=\textbf\def\PY@tc##1{\textcolor[rgb]{0.00,0.50,0.00}{##1}}}
\@namedef{PY@tok@kr}{\let\PY@bf=\textbf\def\PY@tc##1{\textcolor[rgb]{0.00,0.50,0.00}{##1}}}
\@namedef{PY@tok@bp}{\def\PY@tc##1{\textcolor[rgb]{0.00,0.50,0.00}{##1}}}
\@namedef{PY@tok@fm}{\def\PY@tc##1{\textcolor[rgb]{0.00,0.00,1.00}{##1}}}
\@namedef{PY@tok@vc}{\def\PY@tc##1{\textcolor[rgb]{0.10,0.09,0.49}{##1}}}
\@namedef{PY@tok@vg}{\def\PY@tc##1{\textcolor[rgb]{0.10,0.09,0.49}{##1}}}
\@namedef{PY@tok@vi}{\def\PY@tc##1{\textcolor[rgb]{0.10,0.09,0.49}{##1}}}
\@namedef{PY@tok@vm}{\def\PY@tc##1{\textcolor[rgb]{0.10,0.09,0.49}{##1}}}
\@namedef{PY@tok@sa}{\def\PY@tc##1{\textcolor[rgb]{0.73,0.13,0.13}{##1}}}
\@namedef{PY@tok@sb}{\def\PY@tc##1{\textcolor[rgb]{0.73,0.13,0.13}{##1}}}
\@namedef{PY@tok@sc}{\def\PY@tc##1{\textcolor[rgb]{0.73,0.13,0.13}{##1}}}
\@namedef{PY@tok@dl}{\def\PY@tc##1{\textcolor[rgb]{0.73,0.13,0.13}{##1}}}
\@namedef{PY@tok@s2}{\def\PY@tc##1{\textcolor[rgb]{0.73,0.13,0.13}{##1}}}
\@namedef{PY@tok@sh}{\def\PY@tc##1{\textcolor[rgb]{0.73,0.13,0.13}{##1}}}
\@namedef{PY@tok@s1}{\def\PY@tc##1{\textcolor[rgb]{0.73,0.13,0.13}{##1}}}
\@namedef{PY@tok@mb}{\def\PY@tc##1{\textcolor[rgb]{0.40,0.40,0.40}{##1}}}
\@namedef{PY@tok@mf}{\def\PY@tc##1{\textcolor[rgb]{0.40,0.40,0.40}{##1}}}
\@namedef{PY@tok@mh}{\def\PY@tc##1{\textcolor[rgb]{0.40,0.40,0.40}{##1}}}
\@namedef{PY@tok@mi}{\def\PY@tc##1{\textcolor[rgb]{0.40,0.40,0.40}{##1}}}
\@namedef{PY@tok@il}{\def\PY@tc##1{\textcolor[rgb]{0.40,0.40,0.40}{##1}}}
\@namedef{PY@tok@mo}{\def\PY@tc##1{\textcolor[rgb]{0.40,0.40,0.40}{##1}}}
\@namedef{PY@tok@ch}{\let\PY@it=\textit\def\PY@tc##1{\textcolor[rgb]{0.24,0.48,0.48}{##1}}}
\@namedef{PY@tok@cm}{\let\PY@it=\textit\def\PY@tc##1{\textcolor[rgb]{0.24,0.48,0.48}{##1}}}
\@namedef{PY@tok@cpf}{\let\PY@it=\textit\def\PY@tc##1{\textcolor[rgb]{0.24,0.48,0.48}{##1}}}
\@namedef{PY@tok@c1}{\let\PY@it=\textit\def\PY@tc##1{\textcolor[rgb]{0.24,0.48,0.48}{##1}}}
\@namedef{PY@tok@cs}{\let\PY@it=\textit\def\PY@tc##1{\textcolor[rgb]{0.24,0.48,0.48}{##1}}}

\def\PYZbs{\char`\\}
\def\PYZus{\char`\_}
\def\PYZob{\char`\{}
\def\PYZcb{\char`\}}
\def\PYZca{\char`\^}
\def\PYZam{\char`\&}
\def\PYZlt{\char`\<}
\def\PYZgt{\char`\>}
\def\PYZsh{\char`\#}
\def\PYZpc{\char`\%}
\def\PYZdl{\char`\$}
\def\PYZhy{\char`\-}
\def\PYZsq{\char`\'}
\def\PYZdq{\char`\"}
\def\PYZti{\char`\~}
% for compatibility with earlier versions
\def\PYZat{@}
\def\PYZlb{[}
\def\PYZrb{]}
\makeatother


    % For linebreaks inside Verbatim environment from package fancyvrb.
    \makeatletter
        \newbox\Wrappedcontinuationbox
        \newbox\Wrappedvisiblespacebox
        \newcommand*\Wrappedvisiblespace {\textcolor{red}{\textvisiblespace}}
        \newcommand*\Wrappedcontinuationsymbol {\textcolor{red}{\llap{\tiny$\m@th\hookrightarrow$}}}
        \newcommand*\Wrappedcontinuationindent {3ex }
        \newcommand*\Wrappedafterbreak {\kern\Wrappedcontinuationindent\copy\Wrappedcontinuationbox}
        % Take advantage of the already applied Pygments mark-up to insert
        % potential linebreaks for TeX processing.
        %        {, <, #, %, $, ' and ": go to next line.
        %        _, }, ^, &, >, - and ~: stay at end of broken line.
        % Use of \textquotesingle for straight quote.
        \newcommand*\Wrappedbreaksatspecials {%
            \def\PYGZus{\discretionary{\char`\_}{\Wrappedafterbreak}{\char`\_}}%
            \def\PYGZob{\discretionary{}{\Wrappedafterbreak\char`\{}{\char`\{}}%
            \def\PYGZcb{\discretionary{\char`\}}{\Wrappedafterbreak}{\char`\}}}%
            \def\PYGZca{\discretionary{\char`\^}{\Wrappedafterbreak}{\char`\^}}%
            \def\PYGZam{\discretionary{\char`\&}{\Wrappedafterbreak}{\char`\&}}%
            \def\PYGZlt{\discretionary{}{\Wrappedafterbreak\char`\<}{\char`\<}}%
            \def\PYGZgt{\discretionary{\char`\>}{\Wrappedafterbreak}{\char`\>}}%
            \def\PYGZsh{\discretionary{}{\Wrappedafterbreak\char`\#}{\char`\#}}%
            \def\PYGZpc{\discretionary{}{\Wrappedafterbreak\char`\%}{\char`\%}}%
            \def\PYGZdl{\discretionary{}{\Wrappedafterbreak\char`\$}{\char`\$}}%
            \def\PYGZhy{\discretionary{\char`\-}{\Wrappedafterbreak}{\char`\-}}%
            \def\PYGZsq{\discretionary{}{\Wrappedafterbreak\textquotesingle}{\textquotesingle}}%
            \def\PYGZdq{\discretionary{}{\Wrappedafterbreak\char`\"}{\char`\"}}%
            \def\PYGZti{\discretionary{\char`\~}{\Wrappedafterbreak}{\char`\~}}%
        }
        % Some characters . , ; ? ! / are not pygmentized.
        % This macro makes them "active" and they will insert potential linebreaks
        \newcommand*\Wrappedbreaksatpunct {%
            \lccode`\~`\.\lowercase{\def~}{\discretionary{\hbox{\char`\.}}{\Wrappedafterbreak}{\hbox{\char`\.}}}%
            \lccode`\~`\,\lowercase{\def~}{\discretionary{\hbox{\char`\,}}{\Wrappedafterbreak}{\hbox{\char`\,}}}%
            \lccode`\~`\;\lowercase{\def~}{\discretionary{\hbox{\char`\;}}{\Wrappedafterbreak}{\hbox{\char`\;}}}%
            \lccode`\~`\:\lowercase{\def~}{\discretionary{\hbox{\char`\:}}{\Wrappedafterbreak}{\hbox{\char`\:}}}%
            \lccode`\~`\?\lowercase{\def~}{\discretionary{\hbox{\char`\?}}{\Wrappedafterbreak}{\hbox{\char`\?}}}%
            \lccode`\~`\!\lowercase{\def~}{\discretionary{\hbox{\char`\!}}{\Wrappedafterbreak}{\hbox{\char`\!}}}%
            \lccode`\~`\/\lowercase{\def~}{\discretionary{\hbox{\char`\/}}{\Wrappedafterbreak}{\hbox{\char`\/}}}%
            \catcode`\.\active
            \catcode`\,\active
            \catcode`\;\active
            \catcode`\:\active
            \catcode`\?\active
            \catcode`\!\active
            \catcode`\/\active
            \lccode`\~`\~
        }
    \makeatother

    \let\OriginalVerbatim=\Verbatim
    \makeatletter
    \renewcommand{\Verbatim}[1][1]{%
        %\parskip\z@skip
        \sbox\Wrappedcontinuationbox {\Wrappedcontinuationsymbol}%
        \sbox\Wrappedvisiblespacebox {\FV@SetupFont\Wrappedvisiblespace}%
        \def\FancyVerbFormatLine ##1{\hsize\linewidth
            \vtop{\raggedright\hyphenpenalty\z@\exhyphenpenalty\z@
                \doublehyphendemerits\z@\finalhyphendemerits\z@
                \strut ##1\strut}%
        }%
        % If the linebreak is at a space, the latter will be displayed as visible
        % space at end of first line, and a continuation symbol starts next line.
        % Stretch/shrink are however usually zero for typewriter font.
        \def\FV@Space {%
            \nobreak\hskip\z@ plus\fontdimen3\font minus\fontdimen4\font
            \discretionary{\copy\Wrappedvisiblespacebox}{\Wrappedafterbreak}
            {\kern\fontdimen2\font}%
        }%

        % Allow breaks at special characters using \PYG... macros.
        \Wrappedbreaksatspecials
        % Breaks at punctuation characters . , ; ? ! and / need catcode=\active
        \OriginalVerbatim[#1,codes*=\Wrappedbreaksatpunct]%
    }
    \makeatother

    % Exact colors from NB
    \definecolor{incolor}{HTML}{303F9F}
    \definecolor{outcolor}{HTML}{D84315}
    \definecolor{cellborder}{HTML}{CFCFCF}
    \definecolor{cellbackground}{HTML}{F7F7F7}

    % prompt
    \makeatletter
    \newcommand{\boxspacing}{\kern\kvtcb@left@rule\kern\kvtcb@boxsep}
    \makeatother
    \newcommand{\prompt}[4]{
        {\ttfamily\llap{{\color{#2}[#3]:\hspace{3pt}#4}}\vspace{-\baselineskip}}
    }
    

    
    % Prevent overflowing lines due to hard-to-break entities
    \sloppy
    % Setup hyperref package
    \hypersetup{
      breaklinks=true,  % so long urls are correctly broken across lines
      colorlinks=true,
      urlcolor=urlcolor,
      linkcolor=linkcolor,
      citecolor=citecolor,
      }
    % Slightly bigger margins than the latex defaults
    
    \geometry{verbose,tmargin=1in,bmargin=1in,lmargin=1in,rmargin=1in}
    
    

\begin{document}
    
    \maketitle
    
    

    
    \textbf{Studente:} Alessio Camerano\\
\textbf{Matricola:} 1044444\\
\textbf{Corso:} Analisi e Trattamento Segnali Digitali

\begin{center}\rule{0.5\linewidth}{0.5pt}\end{center}

\subsubsection{Indice dei Contenuti}\label{indice-dei-contenuti}

\begin{itemize}
\tightlist
\item
  \textbf{1. Setup del Progetto}

  \begin{itemize}
  \tightlist
  \item
    1.1. Importazione delle librerie
  \item
    1.2. Caricamento dei dati dal sensore
  \item
    1.3. Preprocessing del segnale (rimozione bias e media mobile)
  \end{itemize}
\item
  \textbf{2. Filtraggio del Segnale}

  \begin{itemize}
  \tightlist
  \item
    2.1. Progettazione di un filtro IIR (Butterworth)
  \item
    2.2. Confronto visuale tra segnale grezzo, filtrato FIR e filtrato
    IIR
  \end{itemize}
\item
  \textbf{3. Analisi dei Filtri nel Piano Z}

  \begin{itemize}
  \tightlist
  \item
    3.1. Diagramma poli-zeri del filtro FIR
  \item
    3.2. Diagramma poli-zeri del filtro IIR
  \end{itemize}
\item
  \textbf{4. Analisi nel Dominio del Tempo}

  \begin{itemize}
  \tightlist
  \item
    4.1. Statistiche descrittive
  \item
    4.2. Visualizzazione dei segnali
  \item
    4.3. Analisi di auto-correlazione e cross-correlazione
  \end{itemize}
\item
  \textbf{5. Analisi nel Dominio della Frequenza}

  \begin{itemize}
  \tightlist
  \item
    5.1. Stima della PSD con Periodogramma, Welch e modello AR
  \end{itemize}
\item
  \textbf{6. Analisi Tempo-Frequenza}

  \begin{itemize}
  \tightlist
  \item
    6.1. Spettrogramma (STFT)
  \end{itemize}
\item
  \textbf{7. Conclusioni}

  \begin{itemize}
  \tightlist
  \item
    7.1. Riepilogo dei risultati
  \end{itemize}
\end{itemize}

    \section{1. Setup del Progetto}\label{setup-del-progetto}

    \subsection{1.1. Importazione delle
Librerie}\label{importazione-delle-librerie}

    \begin{tcolorbox}[breakable, size=fbox, boxrule=1pt, pad at break*=1mm,colback=cellbackground, colframe=cellborder]
\begin{Verbatim}[commandchars=\\\{\}]
\PY{k+kn}{import}\PY{+w}{ }\PY{n+nn}{numpy}\PY{+w}{ }\PY{k}{as}\PY{+w}{ }\PY{n+nn}{np}
\PY{k+kn}{import}\PY{+w}{ }\PY{n+nn}{pandas}\PY{+w}{ }\PY{k}{as}\PY{+w}{ }\PY{n+nn}{pd}
\PY{k+kn}{import}\PY{+w}{ }\PY{n+nn}{matplotlib}\PY{n+nn}{.}\PY{n+nn}{pyplot}\PY{+w}{ }\PY{k}{as}\PY{+w}{ }\PY{n+nn}{plt}
\PY{k+kn}{from}\PY{+w}{ }\PY{n+nn}{scipy}\PY{+w}{ }\PY{k+kn}{import} \PY{n}{signal}
\PY{k+kn}{from}\PY{+w}{ }\PY{n+nn}{scipy}\PY{n+nn}{.}\PY{n+nn}{fft}\PY{+w}{ }\PY{k+kn}{import} \PY{n}{fft}\PY{p}{,} \PY{n}{fftfreq}
\PY{k+kn}{from}\PY{+w}{ }\PY{n+nn}{pathlib}\PY{+w}{ }\PY{k+kn}{import} \PY{n}{Path}
\PY{k+kn}{from}\PY{+w}{ }\PY{n+nn}{utils}\PY{+w}{ }\PY{k+kn}{import} \PY{n}{load\PYZus{}gyro\PYZus{}data}\PY{p}{,} \PY{n}{get\PYZus{}sampling\PYZus{}frequency}

\PY{c+c1}{\PYZsh{} \PYZhy{}\PYZhy{}\PYZhy{} Stile e Parametri Globali \PYZhy{}\PYZhy{}\PYZhy{}}
\PY{n}{plt}\PY{o}{.}\PY{n}{style}\PY{o}{.}\PY{n}{use}\PY{p}{(}\PY{l+s+s1}{\PYZsq{}}\PY{l+s+s1}{seaborn\PYZhy{}v0\PYZus{}8\PYZhy{}whitegrid}\PY{l+s+s1}{\PYZsq{}}\PY{p}{)}
\PY{n}{plt}\PY{o}{.}\PY{n}{rcParams}\PY{p}{[}\PY{l+s+s1}{\PYZsq{}}\PY{l+s+s1}{figure.figsize}\PY{l+s+s1}{\PYZsq{}}\PY{p}{]} \PY{o}{=} \PY{p}{(}\PY{l+m+mi}{14}\PY{p}{,} \PY{l+m+mi}{7}\PY{p}{)}
\PY{n}{plt}\PY{o}{.}\PY{n}{rcParams}\PY{p}{[}\PY{l+s+s1}{\PYZsq{}}\PY{l+s+s1}{font.size}\PY{l+s+s1}{\PYZsq{}}\PY{p}{]} \PY{o}{=} \PY{l+m+mi}{10}
\end{Verbatim}
\end{tcolorbox}

    \subsection{1.2. Caricamento dei Dati dal
Sensore}\label{caricamento-dei-dati-dal-sensore}

    \begin{tcolorbox}[breakable, size=fbox, boxrule=1pt, pad at break*=1mm,colback=cellbackground, colframe=cellborder]
\begin{Verbatim}[commandchars=\\\{\}]
\PY{n}{DATA\PYZus{}DIR} \PY{o}{=} \PY{n}{Path}\PY{p}{(}\PY{l+s+s1}{\PYZsq{}}\PY{l+s+s1}{RawData}\PY{l+s+s1}{\PYZsq{}}\PY{p}{)}
\PY{n}{exp\PYZus{}file} \PY{o}{=} \PY{n}{DATA\PYZus{}DIR} \PY{o}{/} \PY{l+s+s1}{\PYZsq{}}\PY{l+s+s1}{gyro\PYZus{}exp12\PYZus{}user06.txt}\PY{l+s+s1}{\PYZsq{}}
\PY{n}{data} \PY{o}{=} \PY{n}{load\PYZus{}gyro\PYZus{}data}\PY{p}{(}\PY{n}{exp\PYZus{}file}\PY{p}{)}
\PY{n}{fs} \PY{o}{=} \PY{n}{get\PYZus{}sampling\PYZus{}frequency}\PY{p}{(}\PY{p}{)}
\PY{n}{T} \PY{o}{=} \PY{l+m+mi}{1} \PY{o}{/} \PY{n}{fs}
\PY{n}{n\PYZus{}samples} \PY{o}{=} \PY{n+nb}{len}\PY{p}{(}\PY{n}{data}\PY{p}{)}
\PY{n}{duration} \PY{o}{=} \PY{n}{n\PYZus{}samples} \PY{o}{*} \PY{n}{T}
\PY{n}{axis\PYZus{}names} \PY{o}{=} \PY{p}{[}\PY{l+s+s1}{\PYZsq{}}\PY{l+s+s1}{gyro\PYZus{}x}\PY{l+s+s1}{\PYZsq{}}\PY{p}{,} \PY{l+s+s1}{\PYZsq{}}\PY{l+s+s1}{gyro\PYZus{}y}\PY{l+s+s1}{\PYZsq{}}\PY{p}{,} \PY{l+s+s1}{\PYZsq{}}\PY{l+s+s1}{gyro\PYZus{}z}\PY{l+s+s1}{\PYZsq{}}\PY{p}{]}
\PY{n+nb}{print}\PY{p}{(}\PY{l+s+sa}{f}\PY{l+s+s2}{\PYZdq{}}\PY{l+s+s2}{Dati caricati da }\PY{l+s+s2}{\PYZsq{}}\PY{l+s+si}{\PYZob{}}\PY{n}{exp\PYZus{}file}\PY{o}{.}\PY{n}{name}\PY{l+s+si}{\PYZcb{}}\PY{l+s+s2}{\PYZsq{}}\PY{l+s+s2}{: }\PY{l+s+si}{\PYZob{}}\PY{n}{n\PYZus{}samples}\PY{l+s+si}{\PYZcb{}}\PY{l+s+s2}{ campioni, durata }\PY{l+s+si}{\PYZob{}}\PY{n}{duration}\PY{l+s+si}{:}\PY{l+s+s2}{.2f}\PY{l+s+si}{\PYZcb{}}\PY{l+s+s2}{s, fs=}\PY{l+s+si}{\PYZob{}}\PY{n}{fs}\PY{l+s+si}{\PYZcb{}}\PY{l+s+s2}{Hz}\PY{l+s+s2}{\PYZdq{}}\PY{p}{)}
\end{Verbatim}
\end{tcolorbox}

    \begin{Verbatim}[commandchars=\\\{\}]
Dati caricati da 'gyro\_exp12\_user06.txt': 32089 campioni, durata 641.78s,
fs=50.0Hz
    \end{Verbatim}

    \subsection{1.3. Preprocessing del Segnale (Rimozione Bias e Media
Mobile)}\label{preprocessing-del-segnale-rimozione-bias-e-media-mobile}

    \begin{tcolorbox}[breakable, size=fbox, boxrule=1pt, pad at break*=1mm,colback=cellbackground, colframe=cellborder]
\begin{Verbatim}[commandchars=\\\{\}]
\PY{k}{def}\PY{+w}{ }\PY{n+nf}{moving\PYZus{}average}\PY{p}{(}\PY{n}{signal\PYZus{}data}\PY{p}{,} \PY{n}{window}\PY{o}{=}\PY{l+m+mi}{5}\PY{p}{)}\PY{p}{:}
\PY{+w}{    }\PY{l+s+sd}{\PYZdq{}\PYZdq{}\PYZdq{}Filtro a media mobile per smussare il segnale.\PYZdq{}\PYZdq{}\PYZdq{}}
    \PY{k}{return} \PY{n}{np}\PY{o}{.}\PY{n}{convolve}\PY{p}{(}\PY{n}{signal\PYZus{}data}\PY{p}{,} \PY{n}{np}\PY{o}{.}\PY{n}{ones}\PY{p}{(}\PY{n}{window}\PY{p}{)}\PY{o}{/}\PY{n}{window}\PY{p}{,} \PY{n}{mode}\PY{o}{=}\PY{l+s+s1}{\PYZsq{}}\PY{l+s+s1}{same}\PY{l+s+s1}{\PYZsq{}}\PY{p}{)}

\PY{k}{def}\PY{+w}{ }\PY{n+nf}{remove\PYZus{}bias}\PY{p}{(}\PY{n}{signal\PYZus{}data}\PY{p}{)}\PY{p}{:}
\PY{+w}{    }\PY{l+s+sd}{\PYZdq{}\PYZdq{}\PYZdq{}Rimuove la componente DC (media) dal segnale.\PYZdq{}\PYZdq{}\PYZdq{}}
    \PY{k}{return} \PY{n}{signal\PYZus{}data} \PY{o}{\PYZhy{}} \PY{n}{np}\PY{o}{.}\PY{n}{mean}\PY{p}{(}\PY{n}{signal\PYZus{}data}\PY{p}{)}

\PY{c+c1}{\PYZsh{} \PYZhy{}\PYZhy{}\PYZhy{} Applicazione Preprocessing \PYZhy{}\PYZhy{}\PYZhy{}}
\PY{n}{data\PYZus{}filt} \PY{o}{=} \PY{n}{pd}\PY{o}{.}\PY{n}{DataFrame}\PY{p}{(}\PY{p}{\PYZob{}}
    \PY{n}{axis}\PY{p}{:} \PY{n}{remove\PYZus{}bias}\PY{p}{(}\PY{n}{moving\PYZus{}average}\PY{p}{(}\PY{n}{data}\PY{p}{[}\PY{n}{axis}\PY{p}{]}\PY{o}{.}\PY{n}{values}\PY{p}{)}\PY{p}{)} \PY{k}{for} \PY{n}{axis} \PY{o+ow}{in} \PY{n}{axis\PYZus{}names}
\PY{p}{\PYZcb{}}\PY{p}{)}
\PY{n+nb}{print}\PY{p}{(}\PY{l+s+s2}{\PYZdq{}}\PY{l+s+s2}{✓ Preprocessing completato: filtro media mobile e rimozione bias applicati.}\PY{l+s+s2}{\PYZdq{}}\PY{p}{)}
\end{Verbatim}
\end{tcolorbox}

    \begin{Verbatim}[commandchars=\\\{\}]
✓ Preprocessing completato: filtro media mobile e rimozione bias applicati.
    \end{Verbatim}

    \section{2. Filtraggio del Segnale}\label{filtraggio-del-segnale}

In questa sezione, confrontiamo il filtro a media mobile (FIR) con un
filtro IIR più selettivo, come il \textbf{Butterworth}. L'obiettivo è
mostrare il trade-off tra la semplicità del FIR e l'efficienza del IIR.
- \textbf{Filtro FIR (Media Mobile)}: Semplice e sempre stabile, ma meno
efficace nel separare le frequenze. - \textbf{Filtro IIR (Butterworth)}:
Offre una transizione più netta tra banda passante e banda attenuata, ma
richiede un'attenta progettazione per garantire la stabilità.

    \subsection{2.1. Progettazione di un Filtro IIR
(Butterworth)}\label{progettazione-di-un-filtro-iir-butterworth}

    \begin{tcolorbox}[breakable, size=fbox, boxrule=1pt, pad at break*=1mm,colback=cellbackground, colframe=cellborder]
\begin{Verbatim}[commandchars=\\\{\}]
\PY{k+kn}{from}\PY{+w}{ }\PY{n+nn}{scipy}\PY{n+nn}{.}\PY{n+nn}{signal}\PY{+w}{ }\PY{k+kn}{import} \PY{n}{butter}\PY{p}{,} \PY{n}{lfilter}

\PY{k}{def}\PY{+w}{ }\PY{n+nf}{butter\PYZus{}lowpass\PYZus{}filter}\PY{p}{(}\PY{n}{data}\PY{p}{,} \PY{n}{cutoff}\PY{p}{,} \PY{n}{fs}\PY{p}{,} \PY{n}{order}\PY{o}{=}\PY{l+m+mi}{5}\PY{p}{)}\PY{p}{:}
\PY{+w}{    }\PY{l+s+sd}{\PYZdq{}\PYZdq{}\PYZdq{}Applica un filtro passa\PYZhy{}basso Butterworth.\PYZdq{}\PYZdq{}\PYZdq{}}
    \PY{n}{nyq} \PY{o}{=} \PY{l+m+mf}{0.5} \PY{o}{*} \PY{n}{fs}
    \PY{n}{normal\PYZus{}cutoff} \PY{o}{=} \PY{n}{cutoff} \PY{o}{/} \PY{n}{nyq}
    \PY{n}{b}\PY{p}{,} \PY{n}{a} \PY{o}{=} \PY{n}{butter}\PY{p}{(}\PY{n}{order}\PY{p}{,} \PY{n}{normal\PYZus{}cutoff}\PY{p}{,} \PY{n}{btype}\PY{o}{=}\PY{l+s+s1}{\PYZsq{}}\PY{l+s+s1}{low}\PY{l+s+s1}{\PYZsq{}}\PY{p}{,} \PY{n}{analog}\PY{o}{=}\PY{k+kc}{False}\PY{p}{)}
    \PY{n}{y} \PY{o}{=} \PY{n}{lfilter}\PY{p}{(}\PY{n}{b}\PY{p}{,} \PY{n}{a}\PY{p}{,} \PY{n}{data}\PY{p}{)}
    \PY{k}{return} \PY{n}{y}

\PY{c+c1}{\PYZsh{} \PYZhy{}\PYZhy{}\PYZhy{} Applicazione Filtro IIR \PYZhy{}\PYZhy{}\PYZhy{}}
\PY{n}{cutoff\PYZus{}freq} \PY{o}{=} \PY{l+m+mi}{5}  \PY{c+c1}{\PYZsh{} Hz}
\PY{n}{data\PYZus{}iir} \PY{o}{=} \PY{n}{pd}\PY{o}{.}\PY{n}{DataFrame}\PY{p}{(}\PY{p}{\PYZob{}}
    \PY{n}{axis}\PY{p}{:} \PY{n}{butter\PYZus{}lowpass\PYZus{}filter}\PY{p}{(}\PY{n}{data}\PY{p}{[}\PY{n}{axis}\PY{p}{]}\PY{o}{.}\PY{n}{values}\PY{p}{,} \PY{n}{cutoff\PYZus{}freq}\PY{p}{,} \PY{n}{fs}\PY{p}{)} \PY{k}{for} \PY{n}{axis} \PY{o+ow}{in} \PY{n}{axis\PYZus{}names}
\PY{p}{\PYZcb{}}\PY{p}{)}
\PY{n+nb}{print}\PY{p}{(}\PY{l+s+sa}{f}\PY{l+s+s2}{\PYZdq{}}\PY{l+s+s2}{✓ Filtro IIR Butterworth applicato con frequenza di taglio a }\PY{l+s+si}{\PYZob{}}\PY{n}{cutoff\PYZus{}freq}\PY{l+s+si}{\PYZcb{}}\PY{l+s+s2}{ Hz.}\PY{l+s+s2}{\PYZdq{}}\PY{p}{)}
\end{Verbatim}
\end{tcolorbox}

    \begin{Verbatim}[commandchars=\\\{\}]
✓ Filtro IIR Butterworth applicato con frequenza di taglio a 5 Hz.
    \end{Verbatim}

    \subsection{2.2. Confronto Visuale tra Segnale Grezzo, Filtrato FIR e
Filtrato
IIR}\label{confronto-visuale-tra-segnale-grezzo-filtrato-fir-e-filtrato-iir}

    \begin{tcolorbox}[breakable, size=fbox, boxrule=1pt, pad at break*=1mm,colback=cellbackground, colframe=cellborder]
\begin{Verbatim}[commandchars=\\\{\}]
\PY{n}{plt}\PY{o}{.}\PY{n}{figure}\PY{p}{(}\PY{n}{figsize}\PY{o}{=}\PY{p}{(}\PY{l+m+mi}{14}\PY{p}{,} \PY{l+m+mi}{7}\PY{p}{)}\PY{p}{)}
\PY{n}{axis\PYZus{}to\PYZus{}plot} \PY{o}{=} \PY{l+s+s1}{\PYZsq{}}\PY{l+s+s1}{gyro\PYZus{}y}\PY{l+s+s1}{\PYZsq{}}
\PY{n}{plt}\PY{o}{.}\PY{n}{plot}\PY{p}{(}\PY{n}{data}\PY{p}{[}\PY{n}{axis\PYZus{}to\PYZus{}plot}\PY{p}{]}\PY{o}{.}\PY{n}{values}\PY{p}{[}\PY{p}{:}\PY{l+m+mi}{500}\PY{p}{]}\PY{p}{,} \PY{n}{label}\PY{o}{=}\PY{l+s+s1}{\PYZsq{}}\PY{l+s+s1}{Segnale Grezzo}\PY{l+s+s1}{\PYZsq{}}\PY{p}{,} \PY{n}{alpha}\PY{o}{=}\PY{l+m+mf}{0.4}\PY{p}{,} \PY{n}{linestyle}\PY{o}{=}\PY{l+s+s1}{\PYZsq{}}\PY{l+s+s1}{\PYZhy{}\PYZhy{}}\PY{l+s+s1}{\PYZsq{}}\PY{p}{)}
\PY{n}{plt}\PY{o}{.}\PY{n}{plot}\PY{p}{(}\PY{n}{data\PYZus{}filt}\PY{p}{[}\PY{n}{axis\PYZus{}to\PYZus{}plot}\PY{p}{]}\PY{o}{.}\PY{n}{values}\PY{p}{[}\PY{p}{:}\PY{l+m+mi}{500}\PY{p}{]}\PY{p}{,} \PY{n}{label}\PY{o}{=}\PY{l+s+s1}{\PYZsq{}}\PY{l+s+s1}{Filtro FIR (Media Mobile)}\PY{l+s+s1}{\PYZsq{}}\PY{p}{,} \PY{n}{linewidth}\PY{o}{=}\PY{l+m+mi}{2}\PY{p}{)}
\PY{n}{plt}\PY{o}{.}\PY{n}{plot}\PY{p}{(}\PY{n}{data\PYZus{}iir}\PY{p}{[}\PY{n}{axis\PYZus{}to\PYZus{}plot}\PY{p}{]}\PY{o}{.}\PY{n}{values}\PY{p}{[}\PY{p}{:}\PY{l+m+mi}{500}\PY{p}{]}\PY{p}{,} \PY{n}{label}\PY{o}{=}\PY{l+s+s1}{\PYZsq{}}\PY{l+s+s1}{Filtro IIR (Butterworth)}\PY{l+s+s1}{\PYZsq{}}\PY{p}{,} \PY{n}{linewidth}\PY{o}{=}\PY{l+m+mi}{2}\PY{p}{)}
\PY{n}{plt}\PY{o}{.}\PY{n}{title}\PY{p}{(}\PY{l+s+sa}{f}\PY{l+s+s1}{\PYZsq{}}\PY{l+s+s1}{Confronto Filtri su Asse }\PY{l+s+si}{\PYZob{}}\PY{n}{axis\PYZus{}to\PYZus{}plot}\PY{o}{.}\PY{n}{upper}\PY{p}{(}\PY{p}{)}\PY{l+s+si}{\PYZcb{}}\PY{l+s+s1}{ (primi 500 campioni)}\PY{l+s+s1}{\PYZsq{}}\PY{p}{)}
\PY{n}{plt}\PY{o}{.}\PY{n}{xlabel}\PY{p}{(}\PY{l+s+s1}{\PYZsq{}}\PY{l+s+s1}{Campioni}\PY{l+s+s1}{\PYZsq{}}\PY{p}{)}
\PY{n}{plt}\PY{o}{.}\PY{n}{ylabel}\PY{p}{(}\PY{l+s+s1}{\PYZsq{}}\PY{l+s+s1}{Velocità Angolare (rad/s)}\PY{l+s+s1}{\PYZsq{}}\PY{p}{)}
\PY{n}{plt}\PY{o}{.}\PY{n}{legend}\PY{p}{(}\PY{p}{)}
\PY{n}{plt}\PY{o}{.}\PY{n}{grid}\PY{p}{(}\PY{k+kc}{True}\PY{p}{,} \PY{n}{which}\PY{o}{=}\PY{l+s+s1}{\PYZsq{}}\PY{l+s+s1}{both}\PY{l+s+s1}{\PYZsq{}}\PY{p}{,} \PY{n}{linestyle}\PY{o}{=}\PY{l+s+s1}{\PYZsq{}}\PY{l+s+s1}{\PYZhy{}\PYZhy{}}\PY{l+s+s1}{\PYZsq{}}\PY{p}{,} \PY{n}{linewidth}\PY{o}{=}\PY{l+m+mf}{0.5}\PY{p}{)}
\PY{n}{plt}\PY{o}{.}\PY{n}{show}\PY{p}{(}\PY{p}{)}
\end{Verbatim}
\end{tcolorbox}

    \begin{center}
    \adjustimage{max size={0.9\linewidth}{0.9\paperheight}}{analisi_segnale_files/analisi_segnale_12_0.png}
    \end{center}
    { \hspace*{\fill} \\}
    
    \section{3. Analisi dei Filtri nel Piano
Z}\label{analisi-dei-filtri-nel-piano-z}

L'analisi dei poli e degli zeri di un filtro nel piano Z è cruciale per
comprenderne il comportamento: - \textbf{Stabilità}: Un filtro è stabile
se tutti i suoi poli si trovano \textbf{all'interno del cerchio
unitario}. I filtri FIR sono intrinsecamente stabili perché i loro poli
sono sempre nell'origine. - \textbf{Risposta in Frequenza}: La posizione
degli \textbf{zeri} sul cerchio unitario indica le frequenze che il
filtro sopprime completamente. La vicinanza dei \textbf{poli} al cerchio
unitario, invece, amplifica le frequenze in quella regione.

    \subsection{3.1. Diagramma Poli-Zeri del Filtro
FIR}\label{diagramma-poli-zeri-del-filtro-fir}

    \begin{tcolorbox}[breakable, size=fbox, boxrule=1pt, pad at break*=1mm,colback=cellbackground, colframe=cellborder]
\begin{Verbatim}[commandchars=\\\{\}]
\PY{k+kn}{from}\PY{+w}{ }\PY{n+nn}{scipy}\PY{n+nn}{.}\PY{n+nn}{signal}\PY{+w}{ }\PY{k+kn}{import} \PY{n}{tf2zpk}
\PY{k+kn}{import}\PY{+w}{ }\PY{n+nn}{matplotlib}\PY{n+nn}{.}\PY{n+nn}{patches}\PY{+w}{ }\PY{k}{as}\PY{+w}{ }\PY{n+nn}{mpatches}

\PY{k}{def}\PY{+w}{ }\PY{n+nf}{plot\PYZus{}zplane}\PY{p}{(}\PY{n}{b}\PY{p}{,} \PY{n}{a}\PY{p}{,} \PY{n}{title}\PY{o}{=}\PY{l+s+s2}{\PYZdq{}}\PY{l+s+s2}{Diagramma Poli\PYZhy{}Zeri}\PY{l+s+s2}{\PYZdq{}}\PY{p}{)}\PY{p}{:}
\PY{+w}{    }\PY{l+s+sd}{\PYZdq{}\PYZdq{}\PYZdq{}Visualizza il diagramma poli\PYZhy{}zeri per un filtro nel piano Z.\PYZdq{}\PYZdq{}\PYZdq{}}
    \PY{n}{fig}\PY{p}{,} \PY{n}{ax} \PY{o}{=} \PY{n}{plt}\PY{o}{.}\PY{n}{subplots}\PY{p}{(}\PY{n}{figsize}\PY{o}{=}\PY{p}{(}\PY{l+m+mi}{6}\PY{p}{,} \PY{l+m+mi}{6}\PY{p}{)}\PY{p}{)}
    
    \PY{c+c1}{\PYZsh{} Disegna il cerchio unitario}
    \PY{n}{unit\PYZus{}circle} \PY{o}{=} \PY{n}{mpatches}\PY{o}{.}\PY{n}{Circle}\PY{p}{(}\PY{p}{(}\PY{l+m+mi}{0}\PY{p}{,} \PY{l+m+mi}{0}\PY{p}{)}\PY{p}{,} \PY{l+m+mi}{1}\PY{p}{,} \PY{n}{fill}\PY{o}{=}\PY{k+kc}{False}\PY{p}{,} \PY{n}{color}\PY{o}{=}\PY{l+s+s1}{\PYZsq{}}\PY{l+s+s1}{gray}\PY{l+s+s1}{\PYZsq{}}\PY{p}{,} \PY{n}{ls}\PY{o}{=}\PY{l+s+s1}{\PYZsq{}}\PY{l+s+s1}{\PYZhy{}\PYZhy{}}\PY{l+s+s1}{\PYZsq{}}\PY{p}{,} \PY{n}{lw}\PY{o}{=}\PY{l+m+mi}{1}\PY{p}{)}
    \PY{n}{ax}\PY{o}{.}\PY{n}{add\PYZus{}patch}\PY{p}{(}\PY{n}{unit\PYZus{}circle}\PY{p}{)}
    
    \PY{c+c1}{\PYZsh{} Trova poli e zeri}
    \PY{n}{z}\PY{p}{,} \PY{n}{p}\PY{p}{,} \PY{n}{k} \PY{o}{=} \PY{n}{tf2zpk}\PY{p}{(}\PY{n}{b}\PY{p}{,} \PY{n}{a}\PY{p}{)}
    
    \PY{c+c1}{\PYZsh{} Plotta zeri (o) e poli (x)}
    \PY{n}{ax}\PY{o}{.}\PY{n}{plot}\PY{p}{(}\PY{n}{np}\PY{o}{.}\PY{n}{real}\PY{p}{(}\PY{n}{z}\PY{p}{)}\PY{p}{,} \PY{n}{np}\PY{o}{.}\PY{n}{imag}\PY{p}{(}\PY{n}{z}\PY{p}{)}\PY{p}{,} \PY{l+s+s1}{\PYZsq{}}\PY{l+s+s1}{o}\PY{l+s+s1}{\PYZsq{}}\PY{p}{,} \PY{n}{markersize}\PY{o}{=}\PY{l+m+mi}{8}\PY{p}{,} \PY{n}{fillstyle}\PY{o}{=}\PY{l+s+s1}{\PYZsq{}}\PY{l+s+s1}{none}\PY{l+s+s1}{\PYZsq{}}\PY{p}{,} \PY{n}{label}\PY{o}{=}\PY{l+s+s1}{\PYZsq{}}\PY{l+s+s1}{Zeri}\PY{l+s+s1}{\PYZsq{}}\PY{p}{)}
    \PY{n}{ax}\PY{o}{.}\PY{n}{plot}\PY{p}{(}\PY{n}{np}\PY{o}{.}\PY{n}{real}\PY{p}{(}\PY{n}{p}\PY{p}{)}\PY{p}{,} \PY{n}{np}\PY{o}{.}\PY{n}{imag}\PY{p}{(}\PY{n}{p}\PY{p}{)}\PY{p}{,} \PY{l+s+s1}{\PYZsq{}}\PY{l+s+s1}{x}\PY{l+s+s1}{\PYZsq{}}\PY{p}{,} \PY{n}{markersize}\PY{o}{=}\PY{l+m+mi}{8}\PY{p}{,} \PY{n}{label}\PY{o}{=}\PY{l+s+s1}{\PYZsq{}}\PY{l+s+s1}{Poli}\PY{l+s+s1}{\PYZsq{}}\PY{p}{)}
    
    \PY{n}{ax}\PY{o}{.}\PY{n}{set\PYZus{}title}\PY{p}{(}\PY{n}{title}\PY{p}{)}
    \PY{n}{ax}\PY{o}{.}\PY{n}{set\PYZus{}xlabel}\PY{p}{(}\PY{l+s+s1}{\PYZsq{}}\PY{l+s+s1}{Parte Reale}\PY{l+s+s1}{\PYZsq{}}\PY{p}{)}
    \PY{n}{ax}\PY{o}{.}\PY{n}{set\PYZus{}ylabel}\PY{p}{(}\PY{l+s+s1}{\PYZsq{}}\PY{l+s+s1}{Parte Imaginaria}\PY{l+s+s1}{\PYZsq{}}\PY{p}{)}
    \PY{n}{ax}\PY{o}{.}\PY{n}{grid}\PY{p}{(}\PY{k+kc}{True}\PY{p}{)}
    \PY{n}{ax}\PY{o}{.}\PY{n}{set\PYZus{}aspect}\PY{p}{(}\PY{l+s+s1}{\PYZsq{}}\PY{l+s+s1}{equal}\PY{l+s+s1}{\PYZsq{}}\PY{p}{)}
    \PY{n}{ax}\PY{o}{.}\PY{n}{legend}\PY{p}{(}\PY{p}{)}
    \PY{n}{plt}\PY{o}{.}\PY{n}{show}\PY{p}{(}\PY{p}{)}

\PY{c+c1}{\PYZsh{} \PYZhy{}\PYZhy{}\PYZhy{} Diagramma per Filtro FIR (Media Mobile) \PYZhy{}\PYZhy{}\PYZhy{}}
\PY{n}{window\PYZus{}size} \PY{o}{=} \PY{l+m+mi}{5}
\PY{n}{b\PYZus{}fir} \PY{o}{=} \PY{n}{np}\PY{o}{.}\PY{n}{ones}\PY{p}{(}\PY{n}{window\PYZus{}size}\PY{p}{)} \PY{o}{/} \PY{n}{window\PYZus{}size}
\PY{n}{a\PYZus{}fir} \PY{o}{=} \PY{n}{np}\PY{o}{.}\PY{n}{array}\PY{p}{(}\PY{p}{[}\PY{l+m+mf}{1.0}\PY{p}{]}\PY{p}{)}
\PY{n}{plot\PYZus{}zplane}\PY{p}{(}\PY{n}{b\PYZus{}fir}\PY{p}{,} \PY{n}{a\PYZus{}fir}\PY{p}{,} \PY{n}{title}\PY{o}{=}\PY{l+s+s2}{\PYZdq{}}\PY{l+s+s2}{Piano Z \PYZhy{} Filtro FIR (Media Mobile)}\PY{l+s+s2}{\PYZdq{}}\PY{p}{)}
\end{Verbatim}
\end{tcolorbox}

    \begin{center}
    \adjustimage{max size={0.9\linewidth}{0.9\paperheight}}{analisi_segnale_files/analisi_segnale_15_0.png}
    \end{center}
    { \hspace*{\fill} \\}
    
    \subsection{3.2. Diagramma Poli-Zeri del Filtro
IIR}\label{diagramma-poli-zeri-del-filtro-iir}

    \begin{tcolorbox}[breakable, size=fbox, boxrule=1pt, pad at break*=1mm,colback=cellbackground, colframe=cellborder]
\begin{Verbatim}[commandchars=\\\{\}]
\PY{n}{nyq} \PY{o}{=} \PY{l+m+mf}{0.5} \PY{o}{*} \PY{n}{fs}
\PY{n}{normal\PYZus{}cutoff} \PY{o}{=} \PY{n}{cutoff\PYZus{}freq} \PY{o}{/} \PY{n}{nyq}
\PY{n}{b\PYZus{}iir}\PY{p}{,} \PY{n}{a\PYZus{}iir} \PY{o}{=} \PY{n}{butter}\PY{p}{(}\PY{l+m+mi}{5}\PY{p}{,} \PY{n}{normal\PYZus{}cutoff}\PY{p}{,} \PY{n}{btype}\PY{o}{=}\PY{l+s+s1}{\PYZsq{}}\PY{l+s+s1}{low}\PY{l+s+s1}{\PYZsq{}}\PY{p}{,} \PY{n}{analog}\PY{o}{=}\PY{k+kc}{False}\PY{p}{)}
\PY{n}{plot\PYZus{}zplane}\PY{p}{(}\PY{n}{b\PYZus{}iir}\PY{p}{,} \PY{n}{a\PYZus{}iir}\PY{p}{,} \PY{n}{title}\PY{o}{=}\PY{l+s+s2}{\PYZdq{}}\PY{l+s+s2}{Piano Z \PYZhy{} Filtro IIR (Butterworth)}\PY{l+s+s2}{\PYZdq{}}\PY{p}{)}
\end{Verbatim}
\end{tcolorbox}

    \begin{center}
    \adjustimage{max size={0.9\linewidth}{0.9\paperheight}}{analisi_segnale_files/analisi_segnale_17_0.png}
    \end{center}
    { \hspace*{\fill} \\}
    
    \subsection{4.1. Statistiche Descrittive}\label{statistiche-descrittive}

    \begin{tcolorbox}[breakable, size=fbox, boxrule=1pt, pad at break*=1mm,colback=cellbackground, colframe=cellborder]
\begin{Verbatim}[commandchars=\\\{\}]
\PY{n+nb}{print}\PY{p}{(}\PY{l+s+s2}{\PYZdq{}}\PY{l+s+se}{\PYZbs{}n}\PY{l+s+s2}{\PYZhy{}\PYZhy{}\PYZhy{} Analisi Statistica Confrontata \PYZhy{}\PYZhy{}\PYZhy{}}\PY{l+s+s2}{\PYZdq{}}\PY{p}{)}
\PY{n}{stats} \PY{o}{=} \PY{n}{pd}\PY{o}{.}\PY{n}{DataFrame}\PY{p}{(}\PY{p}{\PYZob{}}
    \PY{l+s+s1}{\PYZsq{}}\PY{l+s+s1}{Media Grezzo}\PY{l+s+s1}{\PYZsq{}}\PY{p}{:} \PY{n}{data}\PY{p}{[}\PY{n}{axis\PYZus{}names}\PY{p}{]}\PY{o}{.}\PY{n}{mean}\PY{p}{(}\PY{p}{)}\PY{p}{,}
    \PY{l+s+s1}{\PYZsq{}}\PY{l+s+s1}{Media Filtrato}\PY{l+s+s1}{\PYZsq{}}\PY{p}{:} \PY{n}{data\PYZus{}filt}\PY{p}{[}\PY{n}{axis\PYZus{}names}\PY{p}{]}\PY{o}{.}\PY{n}{mean}\PY{p}{(}\PY{p}{)}\PY{p}{,}
    \PY{l+s+s1}{\PYZsq{}}\PY{l+s+s1}{Std Dev Grezzo}\PY{l+s+s1}{\PYZsq{}}\PY{p}{:} \PY{n}{data}\PY{p}{[}\PY{n}{axis\PYZus{}names}\PY{p}{]}\PY{o}{.}\PY{n}{std}\PY{p}{(}\PY{p}{)}\PY{p}{,}
    \PY{l+s+s1}{\PYZsq{}}\PY{l+s+s1}{Std Dev Filtrato}\PY{l+s+s1}{\PYZsq{}}\PY{p}{:} \PY{n}{data\PYZus{}filt}\PY{p}{[}\PY{n}{axis\PYZus{}names}\PY{p}{]}\PY{o}{.}\PY{n}{std}\PY{p}{(}\PY{p}{)}
\PY{p}{\PYZcb{}}\PY{p}{)}
\PY{n+nb}{print}\PY{p}{(}\PY{n}{stats}\PY{p}{)}
\end{Verbatim}
\end{tcolorbox}

    \begin{Verbatim}[commandchars=\\\{\}]

--- Analisi Statistica Confrontata ---
        Media Grezzo  Media Filtrato  Std Dev Grezzo  Std Dev Filtrato
gyro\_x      0.040420    0.000000e+00        0.647423          0.557286
gyro\_y     -0.007699    0.000000e+00        0.673691          0.533480
gyro\_z     -0.010118    6.642863e-18        0.540248          0.443787
    \end{Verbatim}

    \subsection{4.2. Visualizzazione dei
Segnali}\label{visualizzazione-dei-segnali}

    \begin{tcolorbox}[breakable, size=fbox, boxrule=1pt, pad at break*=1mm,colback=cellbackground, colframe=cellborder]
\begin{Verbatim}[commandchars=\\\{\}]
\PY{n}{plt}\PY{o}{.}\PY{n}{figure}\PY{p}{(}\PY{n}{figsize}\PY{o}{=}\PY{p}{(}\PY{l+m+mi}{14}\PY{p}{,} \PY{l+m+mi}{6}\PY{p}{)}\PY{p}{)}
\PY{k}{for} \PY{n}{axis} \PY{o+ow}{in} \PY{n}{axis\PYZus{}names}\PY{p}{:}
    \PY{n}{plt}\PY{o}{.}\PY{n}{plot}\PY{p}{(}\PY{n}{data}\PY{p}{[}\PY{n}{axis}\PY{p}{]}\PY{o}{.}\PY{n}{values}\PY{p}{[}\PY{p}{:}\PY{l+m+mi}{500}\PY{p}{]}\PY{p}{,} \PY{n}{label}\PY{o}{=}\PY{l+s+sa}{f}\PY{l+s+s1}{\PYZsq{}}\PY{l+s+si}{\PYZob{}}\PY{n}{axis}\PY{l+s+si}{\PYZcb{}}\PY{l+s+s1}{ (Grezzo)}\PY{l+s+s1}{\PYZsq{}}\PY{p}{,} \PY{n}{alpha}\PY{o}{=}\PY{l+m+mf}{0.5}\PY{p}{)}
    \PY{n}{plt}\PY{o}{.}\PY{n}{plot}\PY{p}{(}\PY{n}{data\PYZus{}filt}\PY{p}{[}\PY{n}{axis}\PY{p}{]}\PY{o}{.}\PY{n}{values}\PY{p}{[}\PY{p}{:}\PY{l+m+mi}{500}\PY{p}{]}\PY{p}{,} \PY{n}{label}\PY{o}{=}\PY{l+s+sa}{f}\PY{l+s+s1}{\PYZsq{}}\PY{l+s+si}{\PYZob{}}\PY{n}{axis}\PY{l+s+si}{\PYZcb{}}\PY{l+s+s1}{ (Filtrato)}\PY{l+s+s1}{\PYZsq{}}\PY{p}{,} \PY{n}{linewidth}\PY{o}{=}\PY{l+m+mi}{2}\PY{p}{)}
\PY{n}{plt}\PY{o}{.}\PY{n}{title}\PY{p}{(}\PY{l+s+s1}{\PYZsq{}}\PY{l+s+s1}{Confronto Segnale Grezzo vs. Filtrato (primi 500 campioni)}\PY{l+s+s1}{\PYZsq{}}\PY{p}{)}
\PY{n}{plt}\PY{o}{.}\PY{n}{xlabel}\PY{p}{(}\PY{l+s+s1}{\PYZsq{}}\PY{l+s+s1}{Campioni}\PY{l+s+s1}{\PYZsq{}}\PY{p}{)}
\PY{n}{plt}\PY{o}{.}\PY{n}{ylabel}\PY{p}{(}\PY{l+s+s1}{\PYZsq{}}\PY{l+s+s1}{Velocità Angolare (rad/s)}\PY{l+s+s1}{\PYZsq{}}\PY{p}{)}
\PY{n}{plt}\PY{o}{.}\PY{n}{legend}\PY{p}{(}\PY{p}{)}
\PY{n}{plt}\PY{o}{.}\PY{n}{grid}\PY{p}{(}\PY{k+kc}{True}\PY{p}{,} \PY{n}{which}\PY{o}{=}\PY{l+s+s1}{\PYZsq{}}\PY{l+s+s1}{both}\PY{l+s+s1}{\PYZsq{}}\PY{p}{,} \PY{n}{linestyle}\PY{o}{=}\PY{l+s+s1}{\PYZsq{}}\PY{l+s+s1}{\PYZhy{}\PYZhy{}}\PY{l+s+s1}{\PYZsq{}}\PY{p}{,} \PY{n}{linewidth}\PY{o}{=}\PY{l+m+mf}{0.5}\PY{p}{)}
\PY{n}{plt}\PY{o}{.}\PY{n}{show}\PY{p}{(}\PY{p}{)}
\end{Verbatim}
\end{tcolorbox}

    \begin{center}
    \adjustimage{max size={0.9\linewidth}{0.9\paperheight}}{analisi_segnale_files/analisi_segnale_21_0.png}
    \end{center}
    { \hspace*{\fill} \\}
    
    \subsection{4.3. Analisi di Auto-correlazione e
Cross-correlazione}\label{analisi-di-auto-correlazione-e-cross-correlazione}

    \begin{tcolorbox}[breakable, size=fbox, boxrule=1pt, pad at break*=1mm,colback=cellbackground, colframe=cellborder]
\begin{Verbatim}[commandchars=\\\{\}]
\PY{k}{def}\PY{+w}{ }\PY{n+nf}{compute\PYZus{}autocorr}\PY{p}{(}\PY{n}{x}\PY{p}{,} \PY{n}{max\PYZus{}lag}\PY{p}{)}\PY{p}{:}
\PY{+w}{    }\PY{l+s+sd}{\PYZdq{}\PYZdq{}\PYZdq{}Calcola l\PYZsq{}auto\PYZhy{}correlazione normalizzata.\PYZdq{}\PYZdq{}\PYZdq{}}
    \PY{n}{x\PYZus{}norm} \PY{o}{=} \PY{p}{(}\PY{n}{x} \PY{o}{\PYZhy{}} \PY{n}{np}\PY{o}{.}\PY{n}{mean}\PY{p}{(}\PY{n}{x}\PY{p}{)}\PY{p}{)} \PY{o}{/} \PY{n}{np}\PY{o}{.}\PY{n}{std}\PY{p}{(}\PY{n}{x}\PY{p}{)}
    \PY{n}{ac} \PY{o}{=} \PY{n}{np}\PY{o}{.}\PY{n}{correlate}\PY{p}{(}\PY{n}{x\PYZus{}norm}\PY{p}{,} \PY{n}{x\PYZus{}norm}\PY{p}{,} \PY{n}{mode}\PY{o}{=}\PY{l+s+s1}{\PYZsq{}}\PY{l+s+s1}{full}\PY{l+s+s1}{\PYZsq{}}\PY{p}{)}
    \PY{n}{mid} \PY{o}{=} \PY{n+nb}{len}\PY{p}{(}\PY{n}{ac}\PY{p}{)} \PY{o}{/}\PY{o}{/} \PY{l+m+mi}{2}
    \PY{k}{return} \PY{n}{ac}\PY{p}{[}\PY{n}{mid}\PY{p}{:}\PY{n}{mid} \PY{o}{+} \PY{n}{max\PYZus{}lag}\PY{p}{]} \PY{o}{/} \PY{n}{ac}\PY{p}{[}\PY{n}{mid}\PY{p}{]}

\PY{k}{def}\PY{+w}{ }\PY{n+nf}{compute\PYZus{}crosscorr}\PY{p}{(}\PY{n}{x}\PY{p}{,} \PY{n}{y}\PY{p}{,} \PY{n}{max\PYZus{}lag}\PY{p}{)}\PY{p}{:}
\PY{+w}{    }\PY{l+s+sd}{\PYZdq{}\PYZdq{}\PYZdq{}Calcola la cross\PYZhy{}correlazione normalizzata.\PYZdq{}\PYZdq{}\PYZdq{}}
    \PY{n}{x\PYZus{}norm} \PY{o}{=} \PY{p}{(}\PY{n}{x} \PY{o}{\PYZhy{}} \PY{n}{np}\PY{o}{.}\PY{n}{mean}\PY{p}{(}\PY{n}{x}\PY{p}{)}\PY{p}{)} \PY{o}{/} \PY{n}{np}\PY{o}{.}\PY{n}{std}\PY{p}{(}\PY{n}{x}\PY{p}{)}
    \PY{n}{y\PYZus{}norm} \PY{o}{=} \PY{p}{(}\PY{n}{y} \PY{o}{\PYZhy{}} \PY{n}{np}\PY{o}{.}\PY{n}{mean}\PY{p}{(}\PY{n}{y}\PY{p}{)}\PY{p}{)} \PY{o}{/} \PY{n}{np}\PY{o}{.}\PY{n}{std}\PY{p}{(}\PY{n}{y}\PY{p}{)}
    \PY{n}{cc} \PY{o}{=} \PY{n}{np}\PY{o}{.}\PY{n}{correlate}\PY{p}{(}\PY{n}{x\PYZus{}norm}\PY{p}{,} \PY{n}{y\PYZus{}norm}\PY{p}{,} \PY{n}{mode}\PY{o}{=}\PY{l+s+s1}{\PYZsq{}}\PY{l+s+s1}{full}\PY{l+s+s1}{\PYZsq{}}\PY{p}{)}
    \PY{n}{mid} \PY{o}{=} \PY{n+nb}{len}\PY{p}{(}\PY{n}{cc}\PY{p}{)} \PY{o}{/}\PY{o}{/} \PY{l+m+mi}{2}
    \PY{k}{return} \PY{n}{cc}\PY{p}{[}\PY{n}{mid} \PY{o}{\PYZhy{}} \PY{n}{max\PYZus{}lag}\PY{p}{:}\PY{n}{mid} \PY{o}{+} \PY{n}{max\PYZus{}lag}\PY{p}{]} \PY{o}{/} \PY{n+nb}{len}\PY{p}{(}\PY{n}{x}\PY{p}{)}

\PY{c+c1}{\PYZsh{} \PYZhy{}\PYZhy{}\PYZhy{} Auto\PYZhy{}correlazione e Cross\PYZhy{}correlazione \PYZhy{}\PYZhy{}\PYZhy{}}
\PY{n}{fig}\PY{p}{,} \PY{n}{axes} \PY{o}{=} \PY{n}{plt}\PY{o}{.}\PY{n}{subplots}\PY{p}{(}\PY{l+m+mi}{1}\PY{p}{,} \PY{l+m+mi}{2}\PY{p}{,} \PY{n}{figsize}\PY{o}{=}\PY{p}{(}\PY{l+m+mi}{14}\PY{p}{,} \PY{l+m+mi}{5}\PY{p}{)}\PY{p}{)}
\PY{n}{max\PYZus{}lag\PYZus{}ac} \PY{o}{=} \PY{l+m+mi}{250}
\PY{n}{max\PYZus{}lag\PYZus{}cc} \PY{o}{=} \PY{l+m+mi}{50}

\PY{c+c1}{\PYZsh{} Auto\PYZhy{}correlazione}
\PY{n}{lags\PYZus{}ac} \PY{o}{=} \PY{n}{np}\PY{o}{.}\PY{n}{arange}\PY{p}{(}\PY{n}{max\PYZus{}lag\PYZus{}ac}\PY{p}{)} \PY{o}{*} \PY{n}{T}
\PY{n}{ac} \PY{o}{=} \PY{n}{compute\PYZus{}autocorr}\PY{p}{(}\PY{n}{data\PYZus{}filt}\PY{p}{[}\PY{l+s+s1}{\PYZsq{}}\PY{l+s+s1}{gyro\PYZus{}y}\PY{l+s+s1}{\PYZsq{}}\PY{p}{]}\PY{o}{.}\PY{n}{values}\PY{p}{,} \PY{n}{max\PYZus{}lag\PYZus{}ac}\PY{p}{)}
\PY{n}{axes}\PY{p}{[}\PY{l+m+mi}{0}\PY{p}{]}\PY{o}{.}\PY{n}{plot}\PY{p}{(}\PY{n}{lags\PYZus{}ac}\PY{p}{,} \PY{n}{ac}\PY{p}{)}
\PY{n}{axes}\PY{p}{[}\PY{l+m+mi}{0}\PY{p}{]}\PY{o}{.}\PY{n}{set\PYZus{}title}\PY{p}{(}\PY{l+s+s1}{\PYZsq{}}\PY{l+s+s1}{Auto\PYZhy{}correlazione (Asse Y)}\PY{l+s+s1}{\PYZsq{}}\PY{p}{)}
\PY{n}{axes}\PY{p}{[}\PY{l+m+mi}{0}\PY{p}{]}\PY{o}{.}\PY{n}{set\PYZus{}xlabel}\PY{p}{(}\PY{l+s+s1}{\PYZsq{}}\PY{l+s+s1}{Ritardo (s)}\PY{l+s+s1}{\PYZsq{}}\PY{p}{)}
\PY{n}{axes}\PY{p}{[}\PY{l+m+mi}{0}\PY{p}{]}\PY{o}{.}\PY{n}{set\PYZus{}ylabel}\PY{p}{(}\PY{l+s+s1}{\PYZsq{}}\PY{l+s+s1}{Correlazione}\PY{l+s+s1}{\PYZsq{}}\PY{p}{)}
\PY{n}{axes}\PY{p}{[}\PY{l+m+mi}{0}\PY{p}{]}\PY{o}{.}\PY{n}{grid}\PY{p}{(}\PY{k+kc}{True}\PY{p}{)}

\PY{c+c1}{\PYZsh{} Cross\PYZhy{}correlazione}
\PY{n}{lags\PYZus{}cc} \PY{o}{=} \PY{n}{np}\PY{o}{.}\PY{n}{arange}\PY{p}{(}\PY{o}{\PYZhy{}}\PY{n}{max\PYZus{}lag\PYZus{}cc}\PY{p}{,} \PY{n}{max\PYZus{}lag\PYZus{}cc}\PY{p}{)} \PY{o}{*} \PY{n}{T}
\PY{n}{cc} \PY{o}{=} \PY{n}{compute\PYZus{}crosscorr}\PY{p}{(}\PY{n}{data\PYZus{}filt}\PY{p}{[}\PY{l+s+s1}{\PYZsq{}}\PY{l+s+s1}{gyro\PYZus{}x}\PY{l+s+s1}{\PYZsq{}}\PY{p}{]}\PY{o}{.}\PY{n}{values}\PY{p}{,} \PY{n}{data\PYZus{}filt}\PY{p}{[}\PY{l+s+s1}{\PYZsq{}}\PY{l+s+s1}{gyro\PYZus{}y}\PY{l+s+s1}{\PYZsq{}}\PY{p}{]}\PY{o}{.}\PY{n}{values}\PY{p}{,} \PY{n}{max\PYZus{}lag\PYZus{}cc}\PY{p}{)}
\PY{n}{axes}\PY{p}{[}\PY{l+m+mi}{1}\PY{p}{]}\PY{o}{.}\PY{n}{plot}\PY{p}{(}\PY{n}{lags\PYZus{}cc}\PY{p}{,} \PY{n}{cc}\PY{p}{)}
\PY{n}{axes}\PY{p}{[}\PY{l+m+mi}{1}\PY{p}{]}\PY{o}{.}\PY{n}{set\PYZus{}title}\PY{p}{(}\PY{l+s+s1}{\PYZsq{}}\PY{l+s+s1}{Cross\PYZhy{}correlazione (X\PYZhy{}Y)}\PY{l+s+s1}{\PYZsq{}}\PY{p}{)}
\PY{n}{axes}\PY{p}{[}\PY{l+m+mi}{1}\PY{p}{]}\PY{o}{.}\PY{n}{set\PYZus{}xlabel}\PY{p}{(}\PY{l+s+s1}{\PYZsq{}}\PY{l+s+s1}{Ritardo (s)}\PY{l+s+s1}{\PYZsq{}}\PY{p}{)}
\PY{n}{axes}\PY{p}{[}\PY{l+m+mi}{1}\PY{p}{]}\PY{o}{.}\PY{n}{grid}\PY{p}{(}\PY{k+kc}{True}\PY{p}{)}

\PY{n}{plt}\PY{o}{.}\PY{n}{tight\PYZus{}layout}\PY{p}{(}\PY{p}{)}
\PY{n}{plt}\PY{o}{.}\PY{n}{show}\PY{p}{(}\PY{p}{)}
\end{Verbatim}
\end{tcolorbox}

    \begin{center}
    \adjustimage{max size={0.9\linewidth}{0.9\paperheight}}{analisi_segnale_files/analisi_segnale_23_0.png}
    \end{center}
    { \hspace*{\fill} \\}
    
    \section{5. Analisi nel Dominio della
Frequenza}\label{analisi-nel-dominio-della-frequenza}

    \subsection{5.1. Stima della PSD con Periodogramma, Welch e Modello
AR}\label{stima-della-psd-con-periodogramma-welch-e-modello-ar}

    \begin{tcolorbox}[breakable, size=fbox, boxrule=1pt, pad at break*=1mm,colback=cellbackground, colframe=cellborder]
\begin{Verbatim}[commandchars=\\\{\}]
\PY{k+kn}{from}\PY{+w}{ }\PY{n+nn}{scipy}\PY{n+nn}{.}\PY{n+nn}{signal}\PY{+w}{ }\PY{k+kn}{import} \PY{n}{welch}
\PY{k+kn}{from}\PY{+w}{ }\PY{n+nn}{statsmodels}\PY{n+nn}{.}\PY{n+nn}{regression}\PY{n+nn}{.}\PY{n+nn}{linear\PYZus{}model}\PY{+w}{ }\PY{k+kn}{import} \PY{n}{yule\PYZus{}walker}

\PY{c+c1}{\PYZsh{} \PYZhy{}\PYZhy{}\PYZhy{} Funzioni per Stima Spettrale \PYZhy{}\PYZhy{}\PYZhy{}}
\PY{k}{def}\PY{+w}{ }\PY{n+nf}{compute\PYZus{}power\PYZus{}spectrum}\PY{p}{(}\PY{n}{signal\PYZus{}data}\PY{p}{,} \PY{n}{fs}\PY{p}{)}\PY{p}{:}
\PY{+w}{    }\PY{l+s+sd}{\PYZdq{}\PYZdq{}\PYZdq{}Calcola lo spettro di potenza tramite DFT.\PYZdq{}\PYZdq{}\PYZdq{}}
    \PY{n}{N} \PY{o}{=} \PY{n+nb}{len}\PY{p}{(}\PY{n}{signal\PYZus{}data}\PY{p}{)}
    \PY{n}{fft\PYZus{}vals} \PY{o}{=} \PY{n}{fft}\PY{p}{(}\PY{n}{signal\PYZus{}data}\PY{p}{)}
    \PY{n}{freqs} \PY{o}{=} \PY{n}{fftfreq}\PY{p}{(}\PY{n}{N}\PY{p}{,} \PY{l+m+mi}{1}\PY{o}{/}\PY{n}{fs}\PY{p}{)}
    \PY{n}{mask} \PY{o}{=} \PY{n}{freqs} \PY{o}{\PYZgt{}}\PY{o}{=} \PY{l+m+mi}{0}
    \PY{n}{power} \PY{o}{=} \PY{p}{(}\PY{n}{np}\PY{o}{.}\PY{n}{abs}\PY{p}{(}\PY{n}{fft\PYZus{}vals}\PY{p}{[}\PY{n}{mask}\PY{p}{]}\PY{p}{)}\PY{o}{*}\PY{o}{*}\PY{l+m+mi}{2}\PY{p}{)} \PY{o}{/} \PY{n}{N} \PY{c+c1}{\PYZsh{} PSD}
    \PY{k}{return} \PY{n}{freqs}\PY{p}{[}\PY{n}{mask}\PY{p}{]}\PY{p}{,} \PY{n}{power}

\PY{k}{def}\PY{+w}{ }\PY{n+nf}{ar\PYZus{}psd}\PY{p}{(}\PY{n}{x}\PY{p}{,} \PY{n}{order}\PY{p}{,} \PY{n}{fs}\PY{p}{)}\PY{p}{:}
\PY{+w}{    }\PY{l+s+sd}{\PYZdq{}\PYZdq{}\PYZdq{}Stima la PSD con un modello AR(p) via equazioni di Yule\PYZhy{}Walker.\PYZdq{}\PYZdq{}\PYZdq{}}
    \PY{n}{rho}\PY{p}{,} \PY{n}{sigma} \PY{o}{=} \PY{n}{yule\PYZus{}walker}\PY{p}{(}\PY{n}{x}\PY{p}{,} \PY{n}{order}\PY{p}{)}
    \PY{c+c1}{\PYZsh{} Calcola la risposta in frequenza del filtro AR}
    \PY{n}{freqs}\PY{p}{,} \PY{n}{h} \PY{o}{=} \PY{n}{signal}\PY{o}{.}\PY{n}{freqz}\PY{p}{(}\PY{l+m+mi}{1}\PY{p}{,} \PY{n}{np}\PY{o}{.}\PY{n}{r\PYZus{}}\PY{p}{[}\PY{l+m+mi}{1}\PY{p}{,} \PY{o}{\PYZhy{}}\PY{n}{rho}\PY{p}{]}\PY{p}{,} \PY{n}{worN}\PY{o}{=}\PY{l+m+mi}{1024}\PY{p}{,} \PY{n}{fs}\PY{o}{=}\PY{n}{fs}\PY{p}{)}
    \PY{c+c1}{\PYZsh{} Calcola la PSD}
    \PY{n}{psd} \PY{o}{=} \PY{n}{sigma}\PY{o}{*}\PY{o}{*}\PY{l+m+mi}{2} \PY{o}{*} \PY{n}{np}\PY{o}{.}\PY{n}{abs}\PY{p}{(}\PY{n}{h}\PY{p}{)}\PY{o}{*}\PY{o}{*}\PY{l+m+mi}{2}
    \PY{k}{return} \PY{n}{freqs}\PY{p}{,} \PY{n}{psd}

\PY{c+c1}{\PYZsh{} \PYZhy{}\PYZhy{}\PYZhy{} Calcolo e Confronto \PYZhy{}\PYZhy{}\PYZhy{}}
\PY{n}{signal\PYZus{}to\PYZus{}analyze} \PY{o}{=} \PY{n}{data\PYZus{}filt}\PY{p}{[}\PY{l+s+s1}{\PYZsq{}}\PY{l+s+s1}{gyro\PYZus{}y}\PY{l+s+s1}{\PYZsq{}}\PY{p}{]}\PY{o}{.}\PY{n}{values}
\PY{n}{plt}\PY{o}{.}\PY{n}{figure}\PY{p}{(}\PY{n}{figsize}\PY{o}{=}\PY{p}{(}\PY{l+m+mi}{14}\PY{p}{,} \PY{l+m+mi}{6}\PY{p}{)}\PY{p}{)}

\PY{c+c1}{\PYZsh{} 1. DFT (Periodogramma)}
\PY{n}{f\PYZus{}dft}\PY{p}{,} \PY{n}{p\PYZus{}dft} \PY{o}{=} \PY{n}{compute\PYZus{}power\PYZus{}spectrum}\PY{p}{(}\PY{n}{signal\PYZus{}to\PYZus{}analyze}\PY{p}{,} \PY{n}{fs}\PY{p}{)}
\PY{n}{plt}\PY{o}{.}\PY{n}{semilogy}\PY{p}{(}\PY{n}{f\PYZus{}dft}\PY{p}{,} \PY{n}{p\PYZus{}dft}\PY{p}{,} \PY{n}{label}\PY{o}{=}\PY{l+s+s1}{\PYZsq{}}\PY{l+s+s1}{DFT (Periodogramma)}\PY{l+s+s1}{\PYZsq{}}\PY{p}{,} \PY{n}{alpha}\PY{o}{=}\PY{l+m+mf}{0.5}\PY{p}{)}

\PY{c+c1}{\PYZsh{} 2. Welch}
\PY{n}{f\PYZus{}welch}\PY{p}{,} \PY{n}{p\PYZus{}welch} \PY{o}{=} \PY{n}{welch}\PY{p}{(}\PY{n}{signal\PYZus{}to\PYZus{}analyze}\PY{p}{,} \PY{n}{fs}\PY{o}{=}\PY{n}{fs}\PY{p}{,} \PY{n}{nperseg}\PY{o}{=}\PY{l+m+mi}{256}\PY{p}{)}
\PY{n}{plt}\PY{o}{.}\PY{n}{semilogy}\PY{p}{(}\PY{n}{f\PYZus{}welch}\PY{p}{,} \PY{n}{p\PYZus{}welch}\PY{p}{,} \PY{n}{label}\PY{o}{=}\PY{l+s+s1}{\PYZsq{}}\PY{l+s+s1}{Welch}\PY{l+s+s1}{\PYZsq{}}\PY{p}{,} \PY{n}{linewidth}\PY{o}{=}\PY{l+m+mi}{2}\PY{p}{)}

\PY{c+c1}{\PYZsh{} 3. AR}
\PY{n}{f\PYZus{}ar}\PY{p}{,} \PY{n}{p\PYZus{}ar} \PY{o}{=} \PY{n}{ar\PYZus{}psd}\PY{p}{(}\PY{n}{signal\PYZus{}to\PYZus{}analyze}\PY{p}{,} \PY{n}{order}\PY{o}{=}\PY{l+m+mi}{30}\PY{p}{,} \PY{n}{fs}\PY{o}{=}\PY{n}{fs}\PY{p}{)}
\PY{n}{plt}\PY{o}{.}\PY{n}{semilogy}\PY{p}{(}\PY{n}{f\PYZus{}ar}\PY{p}{,} \PY{n}{p\PYZus{}ar}\PY{p}{,} \PY{n}{label}\PY{o}{=}\PY{l+s+s1}{\PYZsq{}}\PY{l+s+s1}{AR(30)}\PY{l+s+s1}{\PYZsq{}}\PY{p}{,} \PY{n}{linewidth}\PY{o}{=}\PY{l+m+mi}{2}\PY{p}{,} \PY{n}{linestyle}\PY{o}{=}\PY{l+s+s1}{\PYZsq{}}\PY{l+s+s1}{\PYZhy{}\PYZhy{}}\PY{l+s+s1}{\PYZsq{}}\PY{p}{)}

\PY{c+c1}{\PYZsh{} \PYZhy{}\PYZhy{}\PYZhy{} Grafico \PYZhy{}\PYZhy{}\PYZhy{}}
\PY{n}{plt}\PY{o}{.}\PY{n}{title}\PY{p}{(}\PY{l+s+s1}{\PYZsq{}}\PY{l+s+s1}{Confronto Metodi di Stima Spettrale (Asse Y)}\PY{l+s+s1}{\PYZsq{}}\PY{p}{)}
\PY{n}{plt}\PY{o}{.}\PY{n}{xlabel}\PY{p}{(}\PY{l+s+s1}{\PYZsq{}}\PY{l+s+s1}{Frequenza (Hz)}\PY{l+s+s1}{\PYZsq{}}\PY{p}{)}
\PY{n}{plt}\PY{o}{.}\PY{n}{ylabel}\PY{p}{(}\PY{l+s+s1}{\PYZsq{}}\PY{l+s+s1}{Power Spectral Density (PSD)}\PY{l+s+s1}{\PYZsq{}}\PY{p}{)}
\PY{n}{plt}\PY{o}{.}\PY{n}{xlim}\PY{p}{(}\PY{l+m+mi}{0}\PY{p}{,} \PY{n}{fs} \PY{o}{/} \PY{l+m+mi}{2}\PY{p}{)}
\PY{n}{plt}\PY{o}{.}\PY{n}{ylim}\PY{p}{(}\PY{n}{bottom}\PY{o}{=}\PY{l+m+mf}{1e\PYZhy{}7}\PY{p}{)}
\PY{n}{plt}\PY{o}{.}\PY{n}{legend}\PY{p}{(}\PY{p}{)}
\PY{n}{plt}\PY{o}{.}\PY{n}{grid}\PY{p}{(}\PY{k+kc}{True}\PY{p}{,} \PY{n}{which}\PY{o}{=}\PY{l+s+s1}{\PYZsq{}}\PY{l+s+s1}{both}\PY{l+s+s1}{\PYZsq{}}\PY{p}{)}
\PY{n}{plt}\PY{o}{.}\PY{n}{show}\PY{p}{(}\PY{p}{)}
\end{Verbatim}
\end{tcolorbox}

    \begin{center}
    \adjustimage{max size={0.9\linewidth}{0.9\paperheight}}{analisi_segnale_files/analisi_segnale_26_0.png}
    \end{center}
    { \hspace*{\fill} \\}
    
    \section{6. Analisi Tempo-Frequenza}\label{analisi-tempo-frequenza}

    \subsection{6.1. Spettrogramma (STFT)}\label{spettrogramma-stft}

Lo spettrogramma mostra come il contenuto in frequenza del segnale
cambia nel tempo. È utile per analizzare segnali non stazionari.

    \begin{tcolorbox}[breakable, size=fbox, boxrule=1pt, pad at break*=1mm,colback=cellbackground, colframe=cellborder]
\begin{Verbatim}[commandchars=\\\{\}]
\PY{k+kn}{from}\PY{+w}{ }\PY{n+nn}{scipy}\PY{n+nn}{.}\PY{n+nn}{signal}\PY{+w}{ }\PY{k+kn}{import} \PY{n}{spectrogram}

\PY{c+c1}{\PYZsh{} \PYZhy{}\PYZhy{}\PYZhy{} Calcolo STFT \PYZhy{}\PYZhy{}\PYZhy{}}
\PY{n}{f\PYZus{}stft}\PY{p}{,} \PY{n}{t\PYZus{}stft}\PY{p}{,} \PY{n}{Sxx} \PY{o}{=} \PY{n}{spectrogram}\PY{p}{(}\PY{n}{data\PYZus{}filt}\PY{p}{[}\PY{l+s+s1}{\PYZsq{}}\PY{l+s+s1}{gyro\PYZus{}y}\PY{l+s+s1}{\PYZsq{}}\PY{p}{]}\PY{o}{.}\PY{n}{values}\PY{p}{,} \PY{n}{fs}\PY{o}{=}\PY{n}{fs}\PY{p}{,} \PY{n}{nperseg}\PY{o}{=}\PY{l+m+mi}{128}\PY{p}{,} \PY{n}{noverlap}\PY{o}{=}\PY{l+m+mi}{96}\PY{p}{)}

\PY{c+c1}{\PYZsh{} \PYZhy{}\PYZhy{}\PYZhy{} Visualizzazione \PYZhy{}\PYZhy{}\PYZhy{}}
\PY{n}{plt}\PY{o}{.}\PY{n}{figure}\PY{p}{(}\PY{n}{figsize}\PY{o}{=}\PY{p}{(}\PY{l+m+mi}{14}\PY{p}{,} \PY{l+m+mi}{6}\PY{p}{)}\PY{p}{)}
\PY{n}{plt}\PY{o}{.}\PY{n}{pcolormesh}\PY{p}{(}\PY{n}{t\PYZus{}stft}\PY{p}{,} \PY{n}{f\PYZus{}stft}\PY{p}{,} \PY{l+m+mi}{10} \PY{o}{*} \PY{n}{np}\PY{o}{.}\PY{n}{log10}\PY{p}{(}\PY{n}{Sxx} \PY{o}{+} \PY{l+m+mf}{1e\PYZhy{}9}\PY{p}{)}\PY{p}{,} \PY{n}{shading}\PY{o}{=}\PY{l+s+s1}{\PYZsq{}}\PY{l+s+s1}{gouraud}\PY{l+s+s1}{\PYZsq{}}\PY{p}{,} \PY{n}{cmap}\PY{o}{=}\PY{l+s+s1}{\PYZsq{}}\PY{l+s+s1}{viridis}\PY{l+s+s1}{\PYZsq{}}\PY{p}{)}
\PY{n}{plt}\PY{o}{.}\PY{n}{colorbar}\PY{p}{(}\PY{n}{label}\PY{o}{=}\PY{l+s+s1}{\PYZsq{}}\PY{l+s+s1}{PSD [dB]}\PY{l+s+s1}{\PYZsq{}}\PY{p}{)}
\PY{n}{plt}\PY{o}{.}\PY{n}{title}\PY{p}{(}\PY{l+s+s1}{\PYZsq{}}\PY{l+s+s1}{Spettrogramma (STFT) \PYZhy{} Asse Y}\PY{l+s+s1}{\PYZsq{}}\PY{p}{)}
\PY{n}{plt}\PY{o}{.}\PY{n}{ylabel}\PY{p}{(}\PY{l+s+s1}{\PYZsq{}}\PY{l+s+s1}{Frequenza (Hz)}\PY{l+s+s1}{\PYZsq{}}\PY{p}{)}
\PY{n}{plt}\PY{o}{.}\PY{n}{xlabel}\PY{p}{(}\PY{l+s+s1}{\PYZsq{}}\PY{l+s+s1}{Tempo (s)}\PY{l+s+s1}{\PYZsq{}}\PY{p}{)}
\PY{n}{plt}\PY{o}{.}\PY{n}{ylim}\PY{p}{(}\PY{l+m+mi}{0}\PY{p}{,} \PY{n}{fs} \PY{o}{/} \PY{l+m+mi}{2}\PY{p}{)}
\PY{n}{plt}\PY{o}{.}\PY{n}{show}\PY{p}{(}\PY{p}{)}
\end{Verbatim}
\end{tcolorbox}

    \begin{center}
    \adjustimage{max size={0.9\linewidth}{0.9\paperheight}}{analisi_segnale_files/analisi_segnale_29_0.png}
    \end{center}
    { \hspace*{\fill} \\}
    
    \section{7. Conclusioni}\label{conclusioni}

\subsection{7.1. Riepilogo dei Risultati}\label{riepilogo-dei-risultati}

Questo progetto ha mostrato un flusso di lavoro completo per l'analisi
di segnali giroscopici: - \textbf{Pulizia del segnale}: Il preprocessing
si è rivelato fondamentale per rimuovere il rumore e le componenti DC. -
\textbf{Analisi Temporale}: L'auto-correlazione ha identificato la
periodicità del movimento, mentre la cross-correlazione ha mostrato la
sincronia tra gli assi. - \textbf{Analisi Spettrale}: Il metodo di Welch
si è dimostrato un ottimo compromesso tra riduzione della varianza e
risoluzione in frequenza, superiore al semplice periodogramma. Il
modello AR offre una visione parametrica e liscia dello spettro. -
\textbf{Analisi Tempo-Frequenza}: Lo spettrogramma ha confermato la
stabilità delle componenti frequenziali nel tempo, caratteristica di un
movimento periodico.


    % Add a bibliography block to the postdoc
    
    
    
\end{document}
